%% =============================================================================
%%  How to study a book with this notebook (generic: fits any book)
%% =============================================================================
\chapter{How to study this book}

This notebook follows \emph{\livroNome} \parencite{livro} lesson by lesson. It
does not reproduce the text of the publication: it carries only the number,
the title and the pages of each lesson, the list of the Bible texts it cites,
and room for your own notes. Use it with the book open beside you.

The publication itself invites the reader to look up the Bible texts it points
to and to compare what the Bible says with what is being taught. This notebook
organises that work, in the spirit of the Beroeans, who received the message
eagerly and \enquote{examined the Scriptures} to confirm what they heard
(\rb{Ac 17:11}).

\section{The cycle of each lesson}

The route combines three well-established ideas about reading and learning:
survey the whole before reading \parencite{robinson1946}, question the text
while reading \parencite{adler1972}, and recall from memory rather than merely
rereading, spacing the reviews over time \parencite{brown2014}.

\begin{center}
\begin{tikzpicture}[
    passo/.style={circle, minimum size=12mm, inner sep=0pt, font=\EBftitulo\Large, text=white},
    rot/.style={font=\sffamily\scriptsize, text width=25mm, align=center, text=tinta}]
  \foreach \i/\ang/\t/\d in {
      1/90/{See the whole}/{title, points, pictures, summary},
      2/18/{Read and mark}/{the idea of each point},
      3/-54/{Check the texts}/{in context, in more than one translation},
      4/-126/{Say it in your own words}/{without looking},
      5/162/{Review and apply}/{after 1 day, 1 week, 1 month}}{
    \node[passo, fill=destaque] (p\i) at (\ang:2.6cm) {\i};
    \node[rot, anchor=center] at (\ang:4.35cm) {\textbf{\t}\\{\color{suave}\d}};
  }
  \foreach \a/\b in {1/2, 2/3, 3/4, 4/5, 5/1}
    \draw[-{Stealth[length=5pt]}, destaque2, thick, shorten <=2pt, shorten >=2pt]
      (p\a) to[bend left=18] (p\b);
  \node[font=\EBfversal\small, text=destaque2!80!black] at (0,0) {Lesson};
\end{tikzpicture}
\end{center}

\paragraph{1. See the whole.} Before reading, go over the title, the questions
and subheadings, the pictures and the summary. Write in the \enquote{Before
reading} field which question the lesson sets out to answer.

\paragraph{2. Read and mark.} Read point by point. For each one, write down
the main idea in a single sentence. The numbers in the \S\ column of the text
table say at which point of the lesson each text appears.

\paragraph{3. Check the texts.} This is the heart of the notebook. Texts with
the \EBlivtipo{leia} badge are the ones the lesson asks you to read; the
\EBlivtipo{citado} ones appear in brackets or in the body of the text. For
each one, follow the routine below and tick the boxes.

\paragraph{4. Say it in your own words.} Close the book and write the central
idea of the lesson. Then answer the review questions from memory. Recalling
works better than rereading \parencite{brown2014}.

\paragraph{5. Review and apply.} Come back to the lesson a day later, a week
later, and at the review of the part. Record the dates in the lesson's header
and in the progress record. Write down one concrete application.

\section{How to check a Bible text}

\begin{ebtabela}{@{}>{\raggedright}p{.28\linewidth} >{\raggedright\arraybackslash}X@{}}
\EBcab{Question} & \EBcab{How to answer it}\\\midrule
What does the context say? & read the whole chapter (or the paragraph, in the letters); note who is speaking, to whom, when and why\\
Do the translations agree? & compare two or three; where they differ, write down the difference\\
Does it hinge on one word? & open the original with the interlinear (\texttt{ferramentas/interlinear.py}) and look at how the word is used elsewhere\\
What point does the lesson draw from it? & write in one sentence how the lesson uses the text and what reasoning links the two\\
What is left open? & record the doubt in the \enquote{To look into} field\\
\end{ebtabela}

\begin{alerta}[Text and context]
A verse on its own can seem to say more, or less, than it says in context.
Reading the whole chapter is the simplest and most useful habit.
\end{alerta}

\section{How to use this notebook}

\paragraph{Printed or digital.} In \texttt{impresso} mode
(\texttt{config.tex}) each lesson carries ruled lines to write on and boxes to
tick. In \texttt{digital} mode only the notes you type appear.

\paragraph{Notes on the computer.} Create the file \texttt{licoes/07.tex}
(lesson number with two digits). Whatever is in it goes into the lesson, in
place of the blank fields. The same resources as in the exegesis notebook are
available: boxes (\texttt{nota}, \texttt{pergunta}, \texttt{lexico}), text in
Greek or Hebrew (\verb|\gr{...}|, \verb|\hb{...}|) and the interlinear.

\begin{nota}[Example of licoes/07.tex]
\small\ttfamily
\textbackslash begin\{campo\}\{Central idea\}\\
\ \ The answer of the lesson, in my own words.\\
\textbackslash end\{campo\}\\
\textbackslash begin\{campo\}\{Text in focus: Ex 34:6\}\\
\ \ \textbackslash interlinear[modo=leitura]\{dados/biblia/ex034\_06.hbo.tsv\}\\
\textbackslash end\{campo\}
\end{nota}

\paragraph{Indexes.} At the end of the notebook you will find the texts most
cited in the book, an index of every Bible text (in biblical order, with the
page of the lesson) and a chart of the 1189 chapters of the Bible for
recording your reading.

\paragraph{Legend of the text table.}
\S\ = point of the lesson; \EBlivtipo{leia} = the lesson asks you to read it;
\EBlivtipo{citado} = cited or mentioned; \EBcaixa\ = tick when you read the
context and when you compare translations.
